\subsection{2026-09-22 --- State of the model: architecture, features, data inventory}
\label{sec:state}

A consolidated description of what the model currently is, what every feature is made of, and
exactly which data we do and do not use. Written because the project has accumulated fifteen
builders, 299 columns and several architectures, and the question of what is actually left to
exploit now matters more than the next variant.

\subsubsection*{Scores, and which are comparable.}

\begin{table}[htbp]
  \centering
  \caption{Current results. The first block is 4-fold CV \emph{means}; the second is a single
  validation year. Comparing across blocks is what made an earlier summary look contradictory.}
  \begin{tabular}{llrr}
    \toprule
    Model & Protocol & Score & SD \\
    \midrule
    Raw features, joint tree   & 4-fold CV mean & 0.7074 & 0.061 \\
    299 engineered features    & 4-fold CV mean & 0.8144 & 0.030 \\
    \textbf{Pruned top-15}     & 4-fold CV mean & \textbf{0.8157} & 0.016 \\
    \midrule
    Pruned top-15              & 2022 only & 0.8210 & --- \\
    Stage 1 (stride-2 fit)     & 2022 only & 0.8390 & --- \\
    Stacked, uniform           & 2022 only & 0.8391 & --- \\
    Stacked, gated on 2021     & 2022 only & 0.8432 & --- \\
    \bottomrule
  \end{tabular}
\end{table}

The 4-fold mean is held down by 2020 (0.7925); 2019/2021/2022 sit at 0.827/0.822/0.821. On the
2022 fold specifically, gated stacking is $+0.022$ over the pruned model.

\subsubsection*{The current model.}
Stage 1 is a single \texttt{XGBClassifier} with \texttt{multi\_strategy="multi\_output\_tree"},
300 trees, depth 6, lr 0.1, predicting all 14 units jointly from 15 shared features. Stage 2 (when
used) is one depth-4 model per unit taking the unit's own stage-1 prediction as
\texttt{base\_margin}, the other 13 predictions as features, and that unit's hydraulic
neighbourhood. Rows are one per (case, hour): 490{,}896 training rows, 2{,}922 cases $\times$ 168
hours, 14 binary targets per row.

\subsubsection*{The fifteen features, and what each is made of.}

\begin{center}
\small
\begin{tabular}{lll}
\toprule
Feature & Definition & Source data \\
\midrule
\texttt{f15\_units\_justified\_*} & $\mathrm{clip}(\text{spread},0)\cdot p_{\max}/p_{\min}$ & price, water value, generator specs \\
\texttt{f15\_cap\_share\_needed\_*} & spread$^+$ normalised by its case max & price, water value \\
\texttt{f14\_upstream\_spread\_*} & price $-$ mean wv of upstream cascade & price, water value, topology graph \\
\texttt{f07\_price\_minus\_terminal\_wv} & price $-$ wv at case start $+$ 8 days & price, water value (boundary) \\
\texttt{f06\_spread\_sys} & price $-$ mean of all 17 water values & price, water value \\
\texttt{f06\_spread\_*} & price $-$ own reservoir's water value & price, water value, topology \\
\texttt{f06x\_margin\_Haukeli} & spread $\times$ plant capacity & price, water value, $p_{\max}$ \\
\texttt{f10\_breakeven\_hours\_Kjela} & startcost $/$ (spread $\times p_{\max}$) & price, wv, \texttt{startcost\_const} \\
\texttt{f05\_wv\_dev\_Kjela} & own wv $-$ system mean wv & water value, topology \\
\bottomrule
\end{tabular}
\end{center}

\finding{\textbf{Every surviving feature is an economic quantity built from price and water value},
with topology and generator specs used only to decide \emph{which} water value and \emph{which}
capacity. No raw price, storage, inflow, calendar or head column survives pruning. Storage and
inflow enter the model only indirectly, through the water values the optimizer itself derived from
them.}

\subsubsection*{Data inventory.}

\paragraph{Used, and how.}
\begin{itemize}\itemsep1pt
  \item \texttt{Historical\_day\_ahead\_price} --- the price term in every spread. Case-relative
        transforms were built and tested but did not survive pruning.
  \item \texttt{Synthetic\_water\_value} --- the other term in every spread; also system mean,
        per-reservoir deviation, and the day$+$8 terminal value.
  \item \texttt{Historical\_volume} --- built into seasonal percentiles and head; \emph{none
        survived}. Enters only implicitly via water value.
  \item \texttt{Historical\_inflow} --- trailing regimes and case-horizon sums; \emph{none
        survived}.
  \item \texttt{Tokke\_Vinje\_topology.yaml} --- reservoir$\leftrightarrow$plant edges and the full
        cascade walk (which water value belongs to which unit); generator $p_{\min}$, $p_{\max}$,
        \texttt{startcost\_const}; \texttt{vol\_head}, \texttt{outlet\_line} and efficiency curves
        were used for head features that did not survive.
  \item \texttt{Constraint\_min\_volume} / \texttt{Constraint\_min\_flow} --- used only inside
        stage-2 per-unit models, wired per plant through the topology.
\end{itemize}

\paragraph{Available and still unused.}
\begin{itemize}\itemsep1pt
  \item \textbf{Cascade water balance.} \texttt{release} $=$ \texttt{storage}$[d]-$\texttt{storage}$[d{+}1]+$\texttt{inflow}$[d]$
        plus upstream plant discharge. This reconstructs \emph{implied production}, which is the
        physical signature of the very decisions we predict. A naive version (natural inflow only)
        gave weak correlations, but it omitted the upstream-discharge term that matters for
        cascaded reservoirs (Section~\ref{sec:pruning}).
  \item \textbf{River \texttt{up\_head\_flow\_curve}} --- head $\to$ maximum flow, i.e.\ hard
        throughput ceilings per river. 37 river records, unused.
  \item \textbf{Tunnel} \texttt{length}, \texttt{loss\_factor}, endpoint heights. 5 records, unused.
  \item \textbf{Flow-dependent efficiency.} \texttt{turb\_eff\_curves} used only at peak;
        \texttt{gen\_eff\_curve}, \texttt{stopcost\_const}, \texttt{main\_loss},
        \texttt{penstock\_loss} unused.
  \item \textbf{Case overlap.} Consecutive cases share 144 wall-clock hours at 88\% label
        agreement; case-to-case input deltas are unused.
  \item \textbf{Case-curve representation.} PCA / Fourier of the 168h price path, unused.
  \item \textbf{Target reformulation.} Count per plant then merit-order expansion (97.8\%/99.9\%
        accurate, Section~\ref{sec:eda-round-2}) --- never attempted.
\end{itemize}

\paragraph{Not in the dataset.}
Production levels in MW (derivable, see above), the optimizer's dual values, demand or reserve
requirements, and maintenance outages. Whether the provided inflow is the \emph{forecast} SHOP
used or the realised actual is unresolved and would change what inflow features mean.


\subsubsection*{Splitting scheme: measured, not assumed.}
The temporal split has been assumed necessary throughout. Tested directly --- same model, same 15
features, four ways of splitting:

\begin{center}
\begin{tabular}{lrr}
\toprule
Split & micro-AUC & vs.\ temporal \\
\midrule
Random ROW (80/20)              & \textbf{0.9571} & $+0.1224$ \\
Random CASE (80/20)             & \textbf{0.9397} & $+0.1050$ \\
Random CASE $+$ 7-day purge     & 0.7889 & $-0.0458$ \\
Temporal (train $<$2022, validate 2022) & 0.8347 & --- \\
\bottomrule
\end{tabular}
\end{center}

\finding{\textbf{Random splitting inflates the score by $+0.10$ to $+0.12$, landing at 0.957 ---
squarely in the range quoted for strong submissions.} The mechanism is \emph{not} temporal feature
leakage, which the competition permits. It is near-duplicate rows: a case is 168 hours of one week,
so a random row split trains on hour 5 and validates on hour 6 of the same case. Random
\emph{case} splits still leak, because consecutive cases overlap on 144 wall-clock hours at 88\%
label agreement (Section~\ref{sec:eda-round-2}). Any shuffled-validation number on this dataset
should be treated as meaningless.}

\finding{\textbf{The purged run is invalid as measured and is recorded as a failed design.}
Purging $\pm$7 days around 584 randomly chosen validation cases removed \textbf{2{,}230 of 2{,}922}
training cases --- 76\% of the data --- so 0.7889 reflects starvation, not a clean estimate. With
daily rolling cases every case sits within 7 days of roughly fourteen others, making random-fold
purging structurally self-defeating. A valid scheme needs \emph{contiguous} validation blocks with
embargo only at the edges.}

\finding{\textbf{This reframes the gap to 0.99.} The hidden test set is 725 consecutive cases in
2023--24, temporally disjoint from training, so a random-split score would not transfer to the
leaderboard. If scores near 0.99 are genuinely achievable there, the likeliest route is not a
validation artefact but exploiting the \emph{physical consequence} of the decisions: reservoir
storage evolves within each case as a direct result of generation, and the volume series covers
the test period. Reconstructing implied production from the cascade water balance therefore moves
from a speculative idea to the leading hypothesis for the missing $0.15$.}

\subsubsection*{Levers not yet pulled.}
\begin{enumerate}\itemsep1pt
  \item \textbf{Validation scheme.} Everything so far uses temporal splits. Random or purged
        case-level splits would cut estimate variance substantially --- tested separately, because
        the risk here is near-duplicate rows, not temporal features.
  \item \textbf{Target reformulation} --- predict per-plant counts, expand by merit order.
  \item \textbf{Cascade water balance} --- implied production.
  \item \textbf{Asymmetric stage 2} --- correction only where stage 1 is weak, keyed off training
        performance rather than an unstable validation gate (Section~\ref{sec:stacked}).
  \item \textbf{Hyperparameter tuning} --- deliberately untouched. Depth, learning rate, and
        regularisation have never been varied; only tree count, which showed 300 beats 900.
\end{enumerate}
